\chapter{The Kodaira--Parshin counterexample}
\label{chap:counterexample}

% archon:covers MR4736527GeometricLocalSystemsOnVeryGeneralCurvesAndIsomonodromy/Counterexample.lean

In this chapter we construct a variation of Hodge structure on
\(\sM_{g,1}\) whose restriction to each fiber of the forgetful map
\(\sM_{g,1} \to \sM_g\) satisfies the hypotheses of
\cref{lem:positive-Hodge-bundle}, and hence yields a flat vector bundle on a
curve whose isomonodromic deformation to a nearby curve is never semistable.
The variation is produced by the Kodaira--Parshin trick. This refutes the
earlier claimed positive answer to the semistability question of Biswas, Heu,
and Hurtubise.

Throughout, \(\sM_{g,n}\) denotes the analytic moduli stack of smooth proper
curves of genus \(g\) with geometrically connected fibers and \(n\) distinct
marked points.

\section{The center-free-group lemma}

We work in the following setup. Let \(\pi : \sM_{g,2} \to \sM_{g,1}\) be the
natural forgetful map, let \(x \in \sM_{g,2}\) be a point, set
\(\bar{x} := \pi(x) \in \sM_{g,1}\), and let \(C^\circ \subset \sM_{g,2}\)
denote the fiber \(\pi^{-1}(\bar{x})\); it is the complement of a point in a
smooth proper connected curve of genus \(g\). There is a natural short exact
sequence
\[
  1 \to \pi_1(C^\circ, x) \to \pi_1(\sM_{g,2}, x) \to \pi_1(\sM_{g,1}, \bar{x})
  \to 1
\]
associated to \(\sM_{g,2} \to \sM_{g,1}\) with fiber \(C^\circ\), obtained from
the Birman exact sequence for mapping class groups. Let \(G\) be a center-free
finite group equipped with a surjection
\(\gamma : \pi_1(C^\circ, x) \twoheadrightarrow G\) taking the conjugacy class of
the loop around the puncture of \(C^\circ\) to a non-identity conjugacy class of
\(G\). Define
\[
  \Gamma \subset \pi_1(\sM_{g,2}, x)
\]
as the set of \(h \in \pi_1(\sM_{g,2}, x)\) such that there exists
\(\widetilde{\gamma}(h) \in G\) with
\(\gamma(h g h^{-1}) = \widetilde{\gamma}(h)\, \gamma(g)\, \widetilde{\gamma}(h)^{-1}\)
for all \(g \in \pi_1(C^\circ, x)\).

\begin{lemma}[Center-free-group lemma]
  \label{lem:center-free-group}
  \lean{LandesmanLitt.centerFreeGroup}
  % SOURCE: [MR4736527 / arXiv:2202.00039], Lemma (lemma:center-free-group), §5
  % (read from references/MR4736527-local-systems-isomonodromy.tex, lines 2003-2012)
  % SOURCE QUOTE: "Keeping notation as in \autoref{subsubsection:setup},
  %	the map $\gamma$ determines a well-defined surjective homomorphism
  %	\begin{align*}
  %		\widetilde{\gamma}: \Gamma  & \rightarrow G\\
  %		h & \mapsto \widetilde{\gamma}(h).
  %	\end{align*}
  %	Moreover, $\Gamma \subset\pi_1(\mathscr{M}_{g,2},x)$ has finite index."
  \textit{Source: MR4736527 (arXiv:2202.00039), §5, Lemma on the center-free
  group.}

  Keeping the notation above, the map \(\gamma\) determines a well-defined
  surjective homomorphism
  \[
    \widetilde{\gamma} : \Gamma \longrightarrow G, \qquad
    h \longmapsto \widetilde{\gamma}(h).
  \]
  Moreover, \(\Gamma \subset \pi_1(\sM_{g,2}, x)\) has finite index.
\end{lemma}

% SOURCE QUOTE PROOF: "We first claim that $\Gamma$ contains $\pi_1(C^\circ, x)$ and surjects
%	onto $G$. Indeed, for $h\in
% \pi_1(C^\circ, x)$, one may take $\widetilde{\gamma}(h)=\gamma(h)$.
% Therefore, the surjectivity of $\gamma$ implies that $\widetilde{\gamma}$ is
% also surjective.
%
%	Next, we claim that for each $h$, $\widetilde{\gamma}(h)$ is uniquely determined.
%	Indeed, suppose $\widetilde{\gamma}(h)$ may be either
%	$\alpha$ and $\beta$. Then we would have $\alpha \gamma(g) \alpha^{-1} =
%	\beta \gamma(g) \beta^{-1}$.
%	Since $\gamma$ is surjective, as shown above, we find $\alpha
%	\beta^{-1}$ lies in the center of $G$, and therefore is trivial. So
%	$\alpha = \beta$.
%
%	The uniqueness of $\widetilde{\gamma}(h)$ just established shows that $\widetilde{\gamma}$ determines a well-defined map.
% This is moreover a homomorphism by the above established uniqueness, because we
% then obtain $\widetilde{\gamma}(h) \widetilde{\gamma}(h') =
% \widetilde{\gamma}(hh')$.
%
% Finally, we claim $\Gamma$ has finite index in $\pi_1(\mathscr{M}_{g,2}, x)$.
% To see this, observe that there is an action of $\pi_1(\mathscr{M}_{g,2}, x)$
% on the set of surjective homomorphisms $\pi_1(C^{\circ},x) \to G$ sending $\phi:
% \pi_1(C^{\circ},x) \to G$ to the map $\phi^h(g) := \phi(hgh^{-1})$.
% By definition, we have $h \in \Gamma$ if and only if $\gamma^h$ is conjugate to $\gamma$.
% In particular, $\Gamma$ contains the stabilizer of $\gamma$ under the action of
% $\pi_1(\mathscr{M}_{g,2}, x)$.
% But this stabilizer has finite index in
% $\pi_1(\mathscr{M}_{g,2}, x)$ because
% $G$ is finite and $\pi_1(C^\circ, x)$ is finitely generated.
% Therefore,
% there are only
% finitely many
% homomorphisms $\pi_1(C^\circ, x) \to G$,
% and, in particular, finitely many such surjective homomorphisms."
\begin{proof}
  First, \(\Gamma\) contains \(\pi_1(C^\circ, x)\) and \(\widetilde{\gamma}\)
  surjects onto \(G\): for \(h \in \pi_1(C^\circ, x)\) one may take
  \(\widetilde{\gamma}(h) = \gamma(h)\), so the surjectivity of \(\gamma\)
  implies that of \(\widetilde{\gamma}\).

  Next, for each \(h \in \Gamma\) the element \(\widetilde{\gamma}(h)\) is
  uniquely determined. Suppose both \(\alpha\) and \(\beta\) work; then
  \(\alpha\, \gamma(g)\, \alpha^{-1} = \beta\, \gamma(g)\, \beta^{-1}\) for all
  \(g\). Since \(\gamma\) is surjective, \(\alpha \beta^{-1}\) lies in the center
  of \(G\), which is trivial because \(G\) is center-free; hence
  \(\alpha = \beta\). This uniqueness shows \(\widetilde{\gamma}\) is a
  well-defined map, and it is a homomorphism because uniqueness forces
  \(\widetilde{\gamma}(h)\, \widetilde{\gamma}(h') = \widetilde{\gamma}(h h')\).

  Finally, \(\Gamma\) has finite index. There is an action of
  \(\pi_1(\sM_{g,2}, x)\) on the set of surjective homomorphisms
  \(\pi_1(C^\circ, x) \to G\) sending \(\phi\) to \(\phi^h(g) := \phi(h g h^{-1})\),
  and by definition \(h \in \Gamma\) if and only if \(\gamma^h\) is conjugate to
  \(\gamma\). Thus \(\Gamma\) contains the stabilizer of \(\gamma\) under this
  action. That stabilizer has finite index because \(G\) is finite and
  \(\pi_1(C^\circ, x)\) is finitely generated, so there are only finitely many
  homomorphisms \(\pi_1(C^\circ, x) \to G\) at all.
\end{proof}

\section{The Kodaira--Parshin trick}

\begin{proposition}[Kodaira--Parshin trick]
  \label{prop:kodaira-parshin}
  \lean{LandesmanLitt.kodairaParshin}
  \uses{lem:center-free-group}
  % SOURCE: [MR4736527 / arXiv:2202.00039], Proposition (proposition:kodaira-parshin), §5
  % (read from references/MR4736527-local-systems-isomonodromy.tex, lines 1940-1952)
  % SOURCE QUOTE: "Let $Y$ denote a Riemann surface of genus $g\geq 1$ with a point $p \in Y$ and
  %	let $Y^{\circ} := Y - p$. Choose a basepoint $y \in Y^{\circ}$.
  %	Suppose $G$ is a finite center-free group with a surjection
  %	$\phi: \pi_1(Y^{\circ},y) \twoheadrightarrow G$ which sends the loop around the puncture $p \in Y$ to
  %	a non-identity element of $G$.
  %	Then there exists a smooth proper relative dimension $1$ map of
  %	analytic stacks $f: \mathscr{X} \to
  %	\mathscr{M}_{g,1}$ so that the fiber over a geometric point $[(C,p)] \in
  %	\mathscr{M}_{g,1}$ is a finite disjoint union of $G$-covers of $C$
  %	ramified at $p$."
  \textit{Source: MR4736527 (arXiv:2202.00039), §5, Kodaira--Parshin trick.}

  Let \(Y\) be a Riemann surface of genus \(g \geq 1\) with a point \(p \in Y\),
  let \(Y^\circ := Y - p\), and choose a basepoint \(y \in Y^\circ\). Suppose
  \(G\) is a finite center-free group with a surjection
  \(\phi : \pi_1(Y^\circ, y) \twoheadrightarrow G\) which sends the loop around
  the puncture \(p \in Y\) to a non-identity element of \(G\). Then there exists
  a smooth proper relative-dimension-\(1\) map of analytic stacks
  \(f : \mathscr{X} \to \sM_{g,1}\) such that the fiber over a geometric point
  \([(C,p)] \in \sM_{g,1}\) is a finite disjoint union of \(G\)-covers of \(C\)
  ramified at \(p\).
\end{proposition}

% SOURCE QUOTE PROOF: "Let $\widetilde{\Gamma}$ be the kernel of the map $\widetilde{\gamma}$ from
% \autoref{lemma:center-free-group}.
% The subgroup $\widetilde{\Gamma}$ corresponds to a finite \'etale cover
% $\mathscr X^\circ \to \mathscr M_{g,2}$.
% Observe that $\mathscr M_{g,2} \subset \mathscr{M}_{g,1} \times_{\mathscr M_g} \mathscr{M}_{g,1}$
% can be viewed as a dense open substack, and
% let $\mathscr{X}$ be the normalization
% of $\mathscr{M}_{g,1} \times_{\mathscr M_g} \mathscr{M}_{g,1}$ in the function
% field of $\mathscr{X}^\circ,$ forming the following cartesian diagram
% ...
% Restricting the natural map $\mathscr{X}\to\mathscr{M}_{g,1} \times_{\mathscr M_g} \mathscr{M}_{g,1}$ to a fiber $C$ of the universal curve $\mathscr{M}_{g,1} \times_{\mathscr M_g} \mathscr{M}_{g,1} \to \mathscr M_{g,1}$ yields a finite disjoint union of $G$-covers of $C$, ramified only over the tautological marked point of $C$.
% We then take our desired relative curve $f: \mathscr X \to \mathscr{M}_{g,1} \times_{\mathscr M_g}
% \mathscr{M}_{g,1} \to \mathscr M_{g,1}$ as the resulting
% composition."
% SOURCE QUOTE (Birman exact sequence input): "We may obtain this from the Birman exact sequence \cite[Theorem
% 4.6]{farbM:a-primer} for mapping class groups,
% after identifying the fundamental group for $\mathscr M_{g,n}$ with the mapping
% class group of an $n$-times punctured genus $g$ surface.
% (The case $n = 0$ follows from contractibility of the universal cover of
% $\mathscr M_g$ \cite[Theorem 10.6]{farbM:a-primer} with covering group given by
% the mapping class group,
% and the case of general $n$ can be deduced from the Birman exact sequence
% \cite[Theorem 4.6]{farbM:a-primer}.)"
\begin{proof}
  \uses{lem:center-free-group}
  This is a black-box construction in mapping-class-group topology; we record
  the argument and its external inputs but do not reprove them, as no Mathlib
  support is available.

  The setup relies on the short exact sequence
  \(1 \to \pi_1(C^\circ, x) \to \pi_1(\sM_{g,2}, x) \to \pi_1(\sM_{g,1}, \bar{x})
  \to 1\). This is obtained from the Birman exact sequence for mapping class
  groups (Farb--Margalit, \emph{A primer on mapping class groups}, Theorem 4.6),
  after identifying \(\pi_1(\sM_{g,n})\) with the mapping class group of an
  \(n\)-times punctured genus-\(g\) surface; the case \(n = 0\) uses the
  contractibility of the universal cover of \(\sM_g\) (Farb--Margalit, Theorem
  10.6), and the general case is deduced from the Birman exact sequence.

  Given this, let \(\widetilde{\Gamma}\) be the kernel of the homomorphism
  \(\widetilde{\gamma}\) of \cref{lem:center-free-group}. Since \(\Gamma\) has
  finite index, \(\widetilde{\Gamma}\) corresponds to a finite étale cover
  \(\mathscr{X}^\circ \to \sM_{g,2}\). Viewing \(\sM_{g,2}\) as a dense open
  substack of \(\sM_{g,1} \times_{\sM_g} \sM_{g,1}\), let \(\mathscr{X}\) be the
  normalization of \(\sM_{g,1} \times_{\sM_g} \sM_{g,1}\) in the function field of
  \(\mathscr{X}^\circ\). Restricting the natural map
  \(\mathscr{X} \to \sM_{g,1} \times_{\sM_g} \sM_{g,1}\) to a fiber \(C\) of the
  universal curve \(\sM_{g,1} \times_{\sM_g} \sM_{g,1} \to \sM_{g,1}\) yields a
  finite disjoint union of \(G\)-covers of \(C\) ramified only over the
  tautological marked point. The desired map \(f : \mathscr{X} \to \sM_{g,1}\) is
  the composition
  \(\mathscr{X} \to \sM_{g,1} \times_{\sM_g} \sM_{g,1} \to \sM_{g,1}\).
\end{proof}

\section{An explicit group: the symmetric group \texorpdfstring{\(S_3\)}{S3}}

\begin{definition}[The Kodaira--Parshin trick for \(S_3\)]
  \label{ex:kodaira-parshin-S3}
  \lean{LandesmanLitt.kodairaParshin_S3}
  \uses{prop:kodaira-parshin, lem:center-free-group}
  % SOURCE: [MR4736527 / arXiv:2202.00039], Example (example:kodaira-parshin), §5
  % (read from references/MR4736527-local-systems-isomonodromy.tex, lines 2076-2091)
  % SOURCE QUOTE: "As a concrete example of a group $G$ to which
  %	\autoref{proposition:kodaira-parshin} applies, we can take $G = S_3$ to be the
  %	symmetric group on three letters and identify $\pi_1(Y^\circ, y)$
  % with the free group on the generators $a_1, \ldots, a_g, b_1, \ldots, b_g$. The group $\pi_1(Y, y)$
  % is generated by $a_1, \ldots, a_g, b_1, \ldots, b_g$ with the relation
  % $\prod_{i=1}^g \left[ a_i, b_i \right]$.
  % Consider the surjection $\phi: \pi_1(C^\circ, y)\twoheadrightarrow S_3$
  % sending $a_1\mapsto (12), b_1 \mapsto (13)$ and sending $a_i \mapsto
  % \mathrm{id}, b_i
  % \mapsto \mathrm{id}$ for $i > 1$.
  % The loop around the puncture maps to
  % $\phi(\prod_{i=1}^g \left[ a_i, b_i \right]) = (12)(13)(12)^{-1}(13)^{-1} =
  % (123) \neq \mathrm{id}$."
  \textit{Source: MR4736527 (arXiv:2202.00039), §5, Example with \(G = S_3\).}

  The symmetric group \(G = S_3\) on three letters is center-free and satisfies
  the hypotheses of \cref{prop:kodaira-parshin}. Identify
  \(\pi_1(Y^\circ, y)\) with the free group on generators
  \(a_1, \ldots, a_g, b_1, \ldots, b_g\); the group \(\pi_1(Y, y)\) is the
  quotient by the single relation \(\prod_{i=1}^g [a_i, b_i]\). Consider the
  surjection \(\phi : \pi_1(C^\circ, y) \twoheadrightarrow S_3\) sending
  \(a_1 \mapsto (12)\), \(b_1 \mapsto (13)\), and \(a_i, b_i \mapsto \id\) for
  \(i > 1\). Then the loop around the puncture maps to
  \[
    \phi\Bigl(\prod_{i=1}^g [a_i, b_i]\Bigr) = (12)(13)(12)^{-1}(13)^{-1}
    = (123) \neq \id,
  \]
  so \(\phi\) sends the loop around the puncture to a non-identity element, as
  required. Hence \(S_3\) provides an explicit instance of
  \cref{prop:kodaira-parshin}.
\end{definition}

\section{The counterexample}

\begin{theorem}[Counterexample to semistability of isomonodromic deformations]
  \label{thm:counterexample}
  \lean{LandesmanLitt.counterexample}
  \uses{prop:kodaira-parshin, ex:kodaira-parshin-S3, lem:positive-Hodge-bundle,
  cor:unstable, def:complex-variation}
  % SOURCE: [MR4736527 / arXiv:2202.00039], Theorem (theorem:counterexample), §1
  % (read from references/MR4736527-local-systems-isomonodromy.tex, lines 327-336)
  % SOURCE QUOTE: "Let $g\geq 2$ be an integer. There exists a vector bundle with flat connection
  % $(\mathscr{F}, \nabla)$ on $\mathscr{M}_{g,1}$ such that for each fiber $C$ of
  % the forgetful morphism $\mathscr{M}_{g,1} \to \mathscr{M}_g$, the restriction of $(\mathscr{F}, \nabla)$ to $C$
  % \begin{enumerate}
  %     \item has semisimple monodromy and
  %     \item is not semistable.
  % \end{enumerate}"
  \textit{Source: MR4736527 (arXiv:2202.00039), §1, Theorem counterexample.}

  Let \(g \geq 2\) be an integer. There exists a vector bundle with flat
  connection \((\sF, \nabla)\) on \(\sM_{g,1}\) such that for each fiber \(C\) of
  the forgetful morphism \(\sM_{g,1} \to \sM_g\), the restriction
  \((\sF, \nabla)|_C\)
  \begin{enumerate}
    \item has semisimple monodromy, and
    \item is not semistable.
  \end{enumerate}
\end{theorem}

% SOURCE QUOTE PROOF: "Let $f: \mathscr X \to \mathscr{M}_{g,1}$ denote the map from
% \autoref{proposition:kodaira-parshin}.
% Concretely, we can take $G = S_3$ and the map $\phi$ as in
% \autoref{proposition:kodaira-parshin} to be that given in
% \autoref{example:kodaira-parshin}.
% Define the local system $\mathbb V := R^1 f_* \mathbb C$ on $\mathscr M_{g,1}$,
% and define $\mathscr{F}$ to be the vector bundle $\mathbb{V}\otimes
% \mathscr{O}$. Note that $\mathscr{F}$ admits a natural (Gauss-Manin) connection $\on{id}\otimes d$. The local system $\mathbb{V}$ evidently underlies a variation of Hodge structure.
% ...
% We first check that $(\mathscr{F}, \nabla)|_C$ has semisimple monodromy. This is
% true for any flat vector bundle arising from the Gauss-Manin connection on the
% cohomology of a family of smooth proper varieties, by work of Deligne
% \cite[Th\'eor\`eme 4.2.6]{deligne:hodge-ii}.
% We now check that $(\mathscr{F}, \nabla)|_C$ is not semistable.
% By \cite[Theorem
% 4]{cataneseD:answer-to-a-question-by-fujita},
% if $X \to C$ is not isotrivial,
% $f_* \omega_f$ is a destabilizing subsheaf of $\mathscr F$.
% It remains to show $X \to C$ is not isotrivial.
% The fiber of $f|_X$ over a point $x\in C$ is a finite disjoint union of
% finite covers of $C$, branched only over $x$. These fibers must vary in moduli as $x$ varies,
% as there are only finitely many non-constant maps between any
% two curves over of genus at least $2$,
% by de Franchis' theorem."
\begin{proof}
  \uses{prop:kodaira-parshin, ex:kodaira-parshin-S3, lem:positive-Hodge-bundle,
  cor:unstable, def:complex-variation}
  Let \(f : \mathscr{X} \to \sM_{g,1}\) be the map from
  \cref{prop:kodaira-parshin}, applied with the center-free group \(G = S_3\) and
  the surjection \(\phi\) of \cref{ex:kodaira-parshin-S3}. Define the local
  system \(\bV := R^1 f_* \C\) on \(\sM_{g,1}\), and let \(\sF := \bV \otimes \sO\)
  be the associated vector bundle, with its natural Gauss--Manin connection
  \(\nabla = \id \otimes d\). The local system \(\bV\) underlies a polarizable
  complex variation of Hodge structure in the sense of
  \cref{def:complex-variation}.

  Fix a fiber \(C\) of the forgetful morphism \(\sM_{g,1} \to \sM_g\) and set
  \(X := f^{-1}(C) \subset \mathscr{X}\).

  \emph{Semisimplicity.} The restriction \((\sF, \nabla)|_C\) has semisimple
  monodromy because it arises from the Gauss--Manin connection on the cohomology
  of a smooth proper family; this is Deligne's semisimplicity theorem (purity of
  the variation of Hodge structure).

  \emph{Non-semistability.} The fiber of \(f|_X : X \to C\) over a point
  \(x \in C\) is a finite disjoint union of finite covers of \(C\) branched only
  over \(x\). As \(x\) varies these fibers vary in moduli, since by de~Franchis'
  theorem there are only finitely many non-constant maps between two fixed curves
  of genus at least \(2\); hence \(f|_X : X \to C\) is not isotrivial.
  Consequently the period map of the variation is non-constant, so the
  Kodaira--Spencer / Higgs field \(\theta_1\) on the top Hodge piece
  \(F^1 \sF|_C \simeq (f|_X)_* \omega_{f|_X}\) is non-zero. The restriction
  \((\sF, \nabla)|_C\) therefore satisfies the hypotheses of
  \cref{lem:positive-Hodge-bundle}, so \(F^1 \sF|_C\) has positive parabolic
  degree, and by \cref{cor:unstable} the bundle \(\sF|_C\) is not semistable.
  (Equivalently, by Catanese--Dettweiler the destabilizing subsheaf is
  \((f|_X)_* \omega_{f|_X}\).)
\end{proof}

\begin{corollary}[Irreducible counterexample]
  \label{cor:counterexample}
  \lean{LandesmanLitt.counterexample_irreducible}
  \uses{thm:counterexample, def:isomonodromy}
  % SOURCE: [MR4736527 / arXiv:2202.00039], Corollary (corollary:counterexample), §1
  % (read from references/MR4736527-local-systems-isomonodromy.tex, lines 343-354)
  % SOURCE QUOTE: "Let $C$ be a smooth projective curve of genus at least $2$. There exists an irreducible flat
  % vector bundle $(E, \nabla)$ on $C$, whose isomonodromic deformations to
  % a nearby curve are never semistable."
  \textit{Source: MR4736527 (arXiv:2202.00039), §1, Corollary counterexample.}

  Let \(C\) be a smooth projective curve of genus at least \(2\). There exists an
  irreducible flat vector bundle \((E, \nabla)\) on \(C\) whose isomonodromic
  deformations to a nearby curve are never semistable.
\end{corollary}

% SOURCE QUOTE PROOF: "The restriction $(\mathscr{F}, \nabla)|_C$ from \autoref{theorem:counterexample} provides a semisimple flat vector bundle, each of whose flat summands has degree zero; by \autoref{theorem:counterexample}(2), its
% isomonodromic deformation to a nearby curve is never semistable. Hence one of
% the irreducible summands of
% $(\mathscr{F}, \nabla)|_C$
% satisfies the statement of the corollary."
\begin{proof}
  \uses{thm:counterexample, def:isomonodromy}
  The restriction \((\sF, \nabla)|_C\) from \cref{thm:counterexample} is a
  semisimple flat vector bundle, each of whose flat summands has degree zero. By
  \cref{thm:counterexample}(2) its isomonodromic deformation
  (\cref{def:isomonodromy}) to a nearby curve is never semistable. Since
  semistability of the whole would follow from semistability of each summand, one
  of the irreducible summands of \((\sF, \nabla)|_C\) must itself have
  non-semistable isomonodromic deformation, and so satisfies the statement.
\end{proof}
